Scheduling real-time applications in the presence of failures is
challenging.  In many domains, some level of output error is tolerable:
application tasks offer different execution modes, where modes with longer
execution times lead to lower error.  This provides an opportunity for
meeting real-time constraints in the presence of failures by trading off
higher error for faster task executions.  In this context, we study the
problem of scheduling a chain of multi-mode tasks onto a failure-prone
platform, allowing for replication, rollback/recovery, adaptive mode
changes after failures, reuse of persistent intermediate outputs, and
coordination across multiple processing elements. The objective is to
minimize the expected error in the final output produced by a deadline,
with a penalty if no output is produced.  For the case in which the
platform comprises a single processing element, we develop a dynamic
program that computes optimal scheduling decisions for the time-discretized
problem under the assumption that a task is never re-executed after it has
successfully completed once.  We also show how to select an appropriate
discretization resolution in practice.  When the platform comprises
multiple processing elements, we propose various heuristics, several of
which build on the above dynamic program.  We demonstrate and/or evaluate
our dynamic programming approach and heuristics via simulated executions
for a fusion energy application derived from a production use case and for
a large set of synthetic workloads.
